% !PRINTED 551
\weberpage{551}
In the difference $2S_0 - S$ there now occur exactly the same
terms as in $S$; only those in which $x+y$ is odd
receive the negative sign, and we can therefore also write:
\weq{12}{\begin{aligned}
\frac{-2\pi}{\beta}\log f(\omega_1)\,f(\omega_2) &= \mathop{\mathrm{Lim}}\limits_{s=1} \sum \frac{(-1)^{x+y}}{[(x-\alpha y)^2+\beta^2y^2]^s},\\
\frac{-2\pi}{\beta}\log f_1(\omega_1)\,f_1(\omega_2) &= \mathop{\mathrm{Lim}}\limits_{s=1} \sum \frac{(-1)^{x}}{[(x-\alpha y)^2+\beta^2y^2]^s},\\
\frac{-2\pi}{\beta}\log f_2(\omega_1)\,f_2(\omega_2) &= \mathop{\mathrm{Lim}}\limits_{s=1} \sum \frac{(-1)^{y}}{[(x-\alpha y)^2+\beta^2y^2]^s},
\end{aligned}}
in which $x, y$ now assume \gesperrt{all} integral pairs of values with the exception
of $0,0$. These three formulae can be derived from one of them by § 34, (13), (14),
for example by replacing, in the second,
\weq{13}{\begin{array}{rcll}
\omega_1,\ \omega_2 & \text{by} & \omega_1+1, & \omega_2-1,\\
& \text{and by} & \dfrac{-1}{\omega_1}, & \dfrac{-1}{\omega_2},\\
\text{that is} &&&\\
\alpha,\ \beta & \text{by} & \alpha+1, & \beta,\\
& \text{and by} & \dfrac{-\alpha}{\alpha^2+\beta^2}, & \dfrac{\beta}{\alpha^2+\beta^2}\text{.}
\end{array}}
\begin{center}\textbf{\S\,146. A Theorem on the Convergence of Series.}\end{center}

If one wished to interchange the symbols Lim and $\Sigma$ with one another
in the expressions (12), § 145, that is, to put $s=1$ under the
summation sign without more ado, one would not obtain unconditionally
convergent series, and it must therefore be investigated
in what sense these formulae then still remain valid.

In order to carry out this investigation, we will first
derive a general theorem from the theory of series, which is a
sharpening of the theorem Bd.\,II, § 196, 3. Let
\weq{1}{\mu_1 \leqq \mu_2 \leqq \mu_3 \leqq \mu_4 \ldots}
be a sequence of \gesperrt{infinitely many} positive numbers, and let $Z(t)$
denote the number of these numbers that are not greater than $t$.
This $Z(t)$ is to have a finite value for every $t$, from which
it then follows that the $\mu_n$ increase to infinity with $n$. But we
will further assume that
\weq{2}{\frac{Z(t)}{t} = a + \frac{\theta}{\sqrt{t}}}
%CONT>>
